\chapter[Band 17: Ganze Zahlen]{Band 17 auf einen Blick}
\noindent\textbf{Ganze Zahlen}\par\medskip
\noindent\emph{Lesart.}
In Klassenformeln sind \(a,b,c,d\in\mathbb N\), in der Ganzzahlarithmetik
\(z,w,u\in\mathbb Z\). Die Zahlzeichen \(0_{\mathbb Z}=\IntZero\) und
\(1_{\mathbb Z}=\IntOne\) bezeichnen die eingebetteten natürlichen Zahlen.
Die später ausdrücklich registrierten Axiome werden zuvor im Quotientenmodell
bewiesen. Die Übersicht unterscheidet diese beiden Rollen.

\begin{BandUebersicht}
\textbf{Differenzenpaare und ihre Äquivalenz}
\UebFormel{\begin{aligned}
D_{\mathbb Z}&\coloneqq\mathbb N\times\mathbb N,\\
(a,b)\IntPairRel(c,d)&\coloneqq a+d=b+c.
\end{aligned}}
\UebQuellen{\FormulaRefAuto{IntegerPairCarrier}[def];
\FormulaRefAuto{IntPairRelDef}[def]}
&
\textit{Äquivalenzrelation}
\UebFormel{\EqRel{D_{\mathbb Z},\IntPairRel}.}
Reflexivität, Symmetrie und Transitivität folgen aus der natürlichen
Addition und ihrer Kürzbarkeit.
\UebQuellen{\FormulaRefAuto{IntegerPairEquivalence}[thm]}
\\ \addlinespace[.8ex]

\textbf{Quotient und Ganzzahlklassen}
\UebFormel{\begin{aligned}
\mathbb Z&\coloneqq\QuotientSet{D_{\mathbb Z}}{\IntPairRel},\\
\IntClass a b&\coloneqq
\EqClass{(a,b)}{\IntPairRel}{D_{\mathbb Z}}.
\end{aligned}}
\UebQuellen{\FormulaRefAuto{IntegersAsQuotient}[def];
\FormulaRefAuto{IntegerClassDef}[def]}
&
\textit{Repräsentanten und Gleichheitskriterium}
\UebFormel{\begin{aligned}
&\IntClass a b=\IntClass c d\\
&\qquad\leftrightarrow a+d=b+c,\\
&z\in\mathbb Z\vdash
\exists a,b\in\mathbb N\;(z=\IntClass a b).
\end{aligned}}
\UebQuellen{\FormulaRefAuto{IntegerClassEquality}[thm];
\FormulaRefAuto{IntegerPairRepresentative}[thm]}
\\ \addlinespace[.8ex]

\textbf{Natürliche Einbettung und Teilmengenkopie}
\UebFormel{\begin{aligned}
&\IntEmbed\colon\mathbb N\to\mathbb Z,\qquad
\IntEmbed(n)\coloneqq\IntClass n0,\\
&\NatCopyInInt\coloneqq\IntEmbed[\mathbb N].
\end{aligned}}
\UebQuellen{\FormulaRefAuto{NaturalIntegerEmbedding}[def];
\FormulaRefAuto{NaturalCopyInIntegers}[def]}
&
\textit{Injektivität und Verträglichkeit}
Für \(m,n\in\mathbb N\):
\UebFormel{\begin{aligned}
&\IntEmbed\colon\mathbb N\inj\mathbb Z,\\
&\IntEmbed(m+n)=\IntEmbed(m)+\IntEmbed(n),\\
&\IntEmbed(mn)=\IntEmbed(m)\IntEmbed(n).
\end{aligned}}
\UebQuellen{\FormulaRefAuto{NaturalIntegerEmbeddingInjective}[thm];
\FormulaRefAuto{NaturalIntegerEmbeddingAdditive}[thm];
\FormulaRefAuto{NaturalIntegerEmbeddingMultiplicative}[thm]}
\\ \addlinespace[.8ex]

\textbf{Negation, Addition und Subtraktion}
\UebFormel{\begin{aligned}
-\IntClass a b&\coloneqq\IntClass b a,\\
\IntClass a b+\IntClass c d
&\coloneqq\IntClass{a+c}{b+d},\\
z-w&\coloneqq z+(-w).
\end{aligned}}
\UebQuellen{\FormulaRefAuto{IntegerNegationDef}[def];
\FormulaRefAuto{IntegerAdditionDef}[def];
\FormulaRefAuto{IntegerSubtractionDef}[def]}
&
\textit{Repräsentantenunabhängigkeit und additive Gesetze}
\UebFormel{\begin{aligned}
&(z+w)+u=z+(w+u),\\
&z+w=w+z,\qquad z+\IntZero=z,\\
&z+(-z)=\IntZero,\qquad-(-z)=z.
\end{aligned}}
\UebQuellen{\FormulaRefAuto{IntegerAdditionCompatible}[thm];
\FormulaRefAuto{IntAddAssociative}[thm];
\FormulaRefAuto{IntAddCommutative}[thm];
\FormulaRefAuto{IntAddZero}[thm];
\FormulaRefAuto{IntAddInverse}[thm];
\FormulaRefAuto{IntegerDoubleNegation}[thm]}
\\ \addlinespace[.8ex]

\textbf{Multiplikation}
\UebFormel{\begin{aligned}
&\IntClass a b\cdot\IntClass c d\\
&\qquad\coloneqq\IntClass{ac+bd}{ad+bc}.
\end{aligned}}
\UebQuellen{\FormulaRefAuto{IntegerMultiplicationDef}[def]}
&
\textit{Multiplikative Gesetze}
\UebFormel{\begin{aligned}
&(zw)u=z(wu),\qquad zw=wz,\\
&z\IntOne=z,\\
&z(w+u)=zw+zu.
\end{aligned}}
\UebQuellen{\FormulaRefAuto{IntegerMultiplicationCompatible}[thm];
\FormulaRefAuto{IntMulAssociative}[thm];
\FormulaRefAuto{IntMulCommutative}[thm];
\FormulaRefAuto{IntMulOne}[thm];
\FormulaRefAuto{IntMulLeftDistributive}[thm]}
\\ \addlinespace[.8ex]

\textbf{Arithmetische Axiome}
Die axiomatische Zusammenfassung fordert die eben modellierten
Abschluss-, Gruppen-, Einheits- und Distributivgesetze, insbesondere
\UebFormel{\begin{aligned}
&z+\IntZero=z,\qquad z+(-z)=\IntZero,\\
&z(w+u)=zw+zu.
\end{aligned}}
\UebQuellen{\FormulaRefAuto{IntegerArithmeticAxioms}[def];
\FormulaRefAuto{IntegerAddZeroAxiom}[ax];
\FormulaRefAuto{IntegerAddInverseAxiom}[ax];
\FormulaRefAuto{IntegerMulLeftDistributiveAxiom}[ax]}
&
\textit{Folgerungen aus dieser Axiomatik}
\UebFormel{\begin{aligned}
&z\IntZero=\IntZero,\qquad-\IntZero=\IntZero,\\
&z(-w)=-(zw),\\
&(-z)(-w)=zw.
\end{aligned}}
\UebQuellen{\FormulaRefAuto{IntegerMulZeroFromAxioms}[thm];
\FormulaRefAuto{IntegerNegativeZero}[thm];
\FormulaRefAuto{IntegerRightNegativeProduct}[thm];
\FormulaRefAuto{IntegerNegativeTimesNegative}[thm]}
\\ \addlinespace[.8ex]

\textbf{Schritte um eins}
Als Abbildungen \(\mathbb Z\to\mathbb Z\) werden eingeführt:
\UebFormel{\begin{aligned}
&(x\mapsto x+\IntOne)(z)\coloneqq z+\IntOne,\\
&(x\mapsto x-\IntOne)(z)\coloneqq z-\IntOne.
\end{aligned}}
\UebQuellen{\FormulaRefAuto{IntegerSuccessorDef}[def];
\FormulaRefAuto{IntegerPredecessorDef}[def]}
&
\textit{Inverse Schritte und Normalform}
\UebFormel{\begin{aligned}
&(z+\IntOne)-\IntOne=z,\\
&(z-\IntOne)+\IntOne=z,\\
&\exists n\in\mathbb N\;
 \bigl(z=\IntEmbed(n)\lor z=-\IntEmbed(n)\bigr).
\end{aligned}}
Insbesondere ist der Schritt nach rechts bijektiv.
\UebQuellen{\FormulaRefAuto{IntegerPredAfterSucc}[thm];
\FormulaRefAuto{IntegerSuccAfterPred}[thm];
\FormulaRefAuto{IntegerSuccessorBijective}[thm];
\FormulaRefAuto{IntegerSignedNormalForm}[thm]}
\\ \addlinespace[.8ex]

\textbf{Quotientenordnung}
\UebFormel{\begin{aligned}
&\IntClass a b\leq_{\mathbb Z}\IntClass c d
\coloneqq a+d\leq c+b,\\
&z<w\coloneqq z\leq w\land z\neq w.
\end{aligned}}
\UebQuellen{\FormulaRefAuto{IntegerOrderDef}[def];
\FormulaRefAuto{IntegerOrderNotation}[def];
\FormulaRefAuto{IntegerStrictOrderDef}[def]}
&
\textit{Totale und translationsverträgliche Ordnung}
\UebFormel{\begin{aligned}
&\TotOrd{\mathbb Z,\leq_{\mathbb Z}},\\
&z\leq w\leftrightarrow z+u\leq w+u,\\
&z<z+\IntOne.
\end{aligned}}
\UebQuellen{\FormulaRefAuto{IntegerOrderTotal}[thm];
\FormulaRefAuto{IntegerOrderTranslationModel}[thm];
\FormulaRefAuto{IntegerSuccessorPositive}[thm]}
\\ \addlinespace[.8ex]

\textbf{Positive ganze Zahlen}
\UebFormel{\PositiveIntegers\coloneqq
\{b\in\mathbb Z\mid\IntZero<b\}.}
\UebQuellen{\FormulaRefAuto{PositiveIntegersDef}[def]}
&
\textit{Kleinste positive Zahl und Faktorkürzung}
Für \(b\in\PositiveIntegers\) gilt
\UebFormel{\begin{aligned}
&\IntOne\leq b,\\
&z\leq w\leftrightarrow zb\leq wb,\\
&zb=wb\vdash z=w.
\end{aligned}}
\UebQuellen{\FormulaRefAuto{PositiveIntegerOneLowerBound}[thm];
\FormulaRefAuto{IntegerPositiveFactorOrder}[thm];
\FormulaRefAuto{IntegerPositiveFactorCancellation}[thm]}
\\ \addlinespace[.8ex]

\textbf{Zweiseitige Peano-Struktur}
Zur geordneten Arithmetik tritt das Induktionsaxiom mit den Daten
\UebFormel{\begin{aligned}
&P(\IntZero),\\
&\forall z\in\mathbb Z\;(P(z)\to P(z+\IntOne)),\\
&\forall z\in\mathbb Z\;(P(z)\to P(z-\IntOne)).
\end{aligned}}
\UebQuellen{\FormulaRefAuto{IntegerPeanoAxioms}[def];
\FormulaRefAuto{IntegerTwoSidedInductionAxiom}[ax]}
&
\textit{Modellnachweis und Induktion im Quotienten}
Aus den links genannten Daten folgt im konstruierten Modell
\UebFormel{\forall z\in\mathbb Z\;P(z).}
Der Quotient erfüllt die vollständige geordnete Ganzzahlarithmetik
einschließlich der inversen Schritte und dieser Induktion.
\UebQuellen{\FormulaRefAuto{IntegerTwoSidedInductionModel}[thm];
\FormulaRefAuto{IntQuotientModelsIntegerPeano}[thm]}
\\
\end{BandUebersicht}
